\section{Negative equality: exclusion of small fixed points}
\label{sec:local}
Throughout this section $\lambda=-a$, $0<a\le 1/2$, and
$d=a^{-2}\ge4$.  The offspring count is either the constant integer $d$
or $\operatorname{Pois}(d)$.  We consider a permutation-invariant fixed
point of the posterior recursion and use
\[
 x=|\zeta|^2,\qquad y=\Re(\zeta^3),\qquad
 \mu=\E x,\quad \nu=\E x^2,\quad \eta=\E y.
\]
All finite decimals in this section and the next denote exact rational
numbers.  In particular, none of the strict comparisons below is a
floating-point comparison.

\begin{proposition}[Local gap at negative equality]
\label{loc:criterion}\label{loc:gap}
Every nontrivial permutation-invariant fixed point at $da^2=1$ satisfies
$\mu>1/20$.  This includes $a=1/2$ in both offspring models.
\end{proposition}

We first establish the pair inequalities used in the proof.  If a current
message has invariants $(X,y_0)$ and an edge message has invariants
$(z,B_0)$, define $d_x,d_e,d_n$ to be the output means of $x,y,x^2$ minus
$X+z,y_0+B_0,X^2+z^2$, respectively.  The output mean here averages the six
relative label permutations with their likelihood tilt.  Conditional on
the permutation orbits of two independent symmetric marginal messages,
these are exactly the relative orientations of their product experiment.
If the child before the edge has invariants $(w,Y)$, then
\[
 z=a^2w,\qquad B_0=-a^3Y.
\]
Set
\[
 S_A=X^2z,\quad S_B=Xz^2,\quad S=S_A+S_B,\quad
 W=y_0z+XB_0,\quad C=\frac52.
\]
(The letter $W$ in this section denotes only this pair expression.)

\begin{lemma}[Certified pair inequalities]\label{loc:pair}
For any current and child posterior vectors, including vectors on the
boundary of the simplex, and any $0\le a\le1/2$,
\begin{align}
 d_n&\le \frac{13}{2}Xz,\qquad
 -\frac{79}{250}Xz\le d_e\le11Xz,\label{loc:coarsepair}\\
 d_x&\le-Xz-2W+2y_0B_0+\frac{51}{20}S,\label{loc:xpair}\\
 6Xz-6W+y_0B_0-\frac{19}{2}S
 &\le d_e\le6Xz-6W+y_0B_0+S,\label{loc:ypair}\\
 C d_n-2d_e&\le-2Xz+22W-7y_0B_0+26S_A+\frac{111}{2}S_B.
 \label{loc:kpair}
\end{align}
\end{lemma}

\begin{proof}
For $a=1/2$ and strictly positive input vectors,
Section~\ref{cert:pair} proves the inequalities by coefficientwise
nonnegativity of fourteen integer polynomials.  That derivation identifies
each cleared polynomial with a pair inequality, including the likelihood
tilt and the cancellation between the two reflection classes.
We now extend those inequalities to all the inputs in the statement.

For general $a\in[0,1/2]$, replace the child vector $m=1+v$
by $1+2av$, a valid convex mixture with the uniform likelihood vector.
Its coloring edge is $1-av$, so the same certificate gives the desired
edge invariants $z=a^2w$ and $B_0=-a^3Y$.

Finally approximate boundary vectors by positive mixtures with the
uniform vector.  For each relative permutation, write $Z$ and
$m^{\mathrm{out}}$ for the normalizer and output likelihood vector
in Equation~\eqref{fnd:product}.  Its contribution to a tilted moment
is $Zf(m^{\mathrm{out}})/6$, where $f=x,y,x^2$ is bounded and
continuous on the simplex.  This contribution is continuous when
$Z>0$ and tends to zero as $Z\to0$.  Thus the tilted expectations extend
continuously over every zero-normalizer orientation, as do the
polynomial right-hand sides.  Taking limits proves the full closed
simplex statement, including $a=1/2$.
\end{proof}

\begin{proof}[Proof of Proposition~\ref{loc:gap}]
Suppose $0<\mu\le m_0:=1/20$.  We introduce a common notation for
summing partial-product increments.  In the regular model let
$(\mu_j,\nu_j,\eta_j)$ be the moments from $j$ edge branches,
$s=j/d$, $0\le j\le d$, and set
\[
 \operatorname{av} f_j=\frac1d\sum_{j=0}^{d-1}f_j,
 \qquad \alpha=1-\frac1d.
\]
A normalized increment is $d$ times the increment from $j$ to $j+1$.
In the Poisson model the partial product has $\operatorname{Pois}(sd)$
branches, $0\le s\le1$; the subscript $j$ below denotes this partial
product, $\operatorname{av}$ means integration in $s$ over $[0,1]$,
and $\alpha=1$.  Its normalized increment is the derivative in $s$.
Differentiating the Poisson series for the bounded functions
$x,y,x^2$ gives $d$ times the expected single-branch increment;
absolute convergence justifies this differentiation.  In both models
the endpoint is the fixed-point law, the initial moments vanish, and
\begin{equation}\label{loc:averages}
 \operatorname{av}s=\frac\alpha2,\qquad
 \operatorname{av}s^2\le\frac\alpha3.
\end{equation}
For regular counts the second average is
$(d-1)(2d-1)/(6d^2)$; for Poisson it is $1/3$.

At equality the standalone edge contributions to $x$ sum to
$da^2\mu=\mu$, which is already the fixed-point endpoint moment.
The average of the actual normalized pair excess in $x$ is therefore
zero.  We will contradict this identity by showing that the pair upper
bound has negative average when $\mu\le m_0$.  To control its higher
moments, we first bound the partial-product moments and then determine
the sign of $\eta$.

The projection bound gives, respectively,
\[
 \mu_j\ge\frac{j a^2\mu}{1+(j-1)a^2\mu},\qquad
 \mu_j\ge\frac{s\mu}{1+s\mu}.
\]
For completeness, one edge posterior complex mean has correlation
$a^2\mu$ with the spin, conditional mean $a^2\mu$ times that spin,
and second moment $a^2\mu$.  The sum of $j$ independent branches
conditional on the spin consequently has second moment
$ja^2\mu+j(j-1)(a^2\mu)^2$ and correlation $ja^2\mu$.
Conditional expectation followed by Cauchy--Schwarz proves the first
bound.  Poisson factorial moments $sd$ and $(sd)^2$ prove the second.
It follows that
\begin{equation}\label{loc:partialmu}
 (s-m_0s^2)\mu\le\mu_j\le1.14s\mu.
\end{equation}
Here is a derivation of the upper bound.  Every simplex message satisfies
\begin{equation}\label{loc:ratio}
 -\frac{x}{2}\le y\le x.
\end{equation}
For $x>0$ this is Equation~\eqref{fnd:ratio-domain}, and at
$x=0$ both sides vanish by Equation~\eqref{fnd:invariant-bounds}.
Thus, when $Xz>0$,
\[
 -\tfrac12\le y_0/X\le1,\qquad
 -\tfrac12\le B_0/z=-aY/w\le\tfrac14,\qquad z\le\tfrac14.
\]
Checking the four corners of this ratio rectangle gives
$-2(r+t)+2rt\le5/2$.  Equation~\eqref{loc:xpair} therefore gives
\[
 \frac{d_x}{Xz}\le-1+\frac52+\frac{51}{20}(1+1/4)
       =\frac{75}{16}<5.
\]
If $Xz=0$, one input is uniform and the excess vanishes.  Hence the
pair gain of $x$ is at most $z(1+5X)$.  Iterating, or solving the
corresponding differential inequality, gives
$\mu_j\le(e^{5s\mu}-1)/5$.
The function $(e^t-1)/t$ increases for $t\ge0$, and
\[
 e^{1/4}\le1+\frac14+\frac1{32}
       +\frac{(1/4)^3/6}{1-1/16}
       =\frac{1849}{1440}<1.285.
\]
This proves the upper bound in~\eqref{loc:partialmu}.

The standalone normalized increments in $x,y,x^2$ are
$\mu,-a\eta,a^2\nu$, respectively.  The partial product and the next
edge branch are independent under
the product marginal law, and each has a symmetric law.  Conditional
on their two permutation orbits, Lemma~\ref{loc:pair} therefore applies;
averaging its bounds factors every mixed invariant moment.  In particular,
Equation~\eqref{loc:coarsepair} gives
\[
 (1-a^2)\nu\le6.5\mu\operatorname{av}\mu_j,
 \qquad
 -.316\mu\operatorname{av}\mu_j
 \le(1+a)\eta\le11\mu\operatorname{av}\mu_j.
\]
Using Equations~\eqref{loc:partialmu} and~\eqref{loc:averages},
$1-a^2\ge3/4$, $a/(1+a)\le1/3$, and $\alpha\le1$, we obtain
\begin{equation}\label{loc:coarse}
 \frac{\nu}{\mu^2}\le4.94,\quad
 \frac{a^2\nu}{\mu^2}\le1.235,\quad
 -.061\le\frac{a\eta}{\mu^2}\le2.09.
\end{equation}
The endpoint bound for $\nu/\mu^2$ before rounding is
$6.5(1.14)\alpha/[2(1-a^2)]$.  The upper bound for $\eta/\mu^2$ is
$11(1.14)\alpha/[2(1+a)]$, and the bound on its negative magnitude is
$.316(1.14)\alpha/[2(1+a)]$.
Summing to a partial endpoint also yields
\begin{align}
 \frac{\nu_j}{\mu^2}&\le1.235s+3.705s^2,\label{loc:partialnu}\\
 -2.09s-.181s^2
 &\le\frac{\eta_j}{\mu^2}\le .061s+6.27s^2.
 \label{loc:partialeta}
\end{align}
For example the quadratic coefficients come from
$6.5(1.14)/2=3.705$,
$.316(1.14)/2=.18012<.181$, and
$11(1.14)/2=6.27$.
In the discrete case we used
$d^{-1}\sum_{i=0}^{j-1}(i/d)^k\le s^{k+1}/(k+1)$;
in the Poisson case this is the exact integral of the monomial.

The sharper $y$ inequality improves the sign information.
Independence gives
\[
 d\E W=\mu\eta_j-a\eta\mu_j,\qquad
 d\E(y_0B_0)=-a\eta\eta_j,\qquad
 d\E S=\mu\nu_j+a^2\nu\mu_j.
\]
Thus the normalized $y$ excess is bounded below and above by
\begin{equation}\label{loc:refinedincrement}
 (6\mu+6a\eta+c a^2\nu)\mu_j
 +(-6\mu-a\eta)\eta_j+c\mu\nu_j,
 \qquad c=-9.5,\ 1,
\end{equation}
respectively.  The coefficient of $\eta_j$ is negative, since
$6\mu+a\eta\ge(6-.061m_0)\mu>0$.
For the lower estimate the coefficient of $\mu_j$ is positive, since
\[
 6\mu+6a\eta-9.5a^2\nu
 \ge[6-(6(.061)+9.5(1.235))(.05)]\mu
 =5.395075\mu>0.
\]
These observations determine the valid directions of all substitutions.
Adding the standalone increment and averaging, define
$R=(1+a)\eta/(\alpha\mu^2)$.  Equations
\eqref{loc:partialmu}--\eqref{loc:partialeta} give
\begin{align*}
 R&\le6.689(.57)+(6+.105)(.05)(2.09/2+.181/3)
       +.05(1.235/2+3.705/3)\\
  &=\frac{2121379}{500000}<4.25,\\
 R&\ge[6-(6(.061)+9.5(1.235))(.05)](1/2-.05/3)\\
 &\hspace{1em}-(6+.105)(.05)(.061/2+6.27/3)
       -9.5(.05)(1.235/2+3.705/3)\\
  &=\frac{25929587}{24000000}>0.
\end{align*}
Here $6+(6(2.09)+1.235)(.05)=6.68875<6.689$,
$2.09(.05)=.1045<.105$, and $1.14/2=.57$.
Consequently
\begin{equation}\label{loc:etapositive}
 \eta\ge0,\qquad 0\le\frac{a\eta}{\mu^2}
       \le\frac{4.25}{3}<1.417.
\end{equation}
The standalone $y$ increment is now nonpositive.  Keeping the preceding
coarse partial bounds in the upper expression of
\eqref{loc:refinedincrement}, its coefficients of $s$ and $s^2$ after
division by $\mu^2$ are at most
\[
 6.689(1.14)+(6+.105)(.05)(2.09)+.05(1.235)
   =8.3251825<8.34,
\]
\[
 (6+.105)(.05)(.181)+.05(3.705)=.24050025<.242.
\]
Summation or integration therefore improves the partial bound to
\begin{equation}\label{loc:etaimproved}
 \eta_j/\mu^2\le4.17s^2+.081s^3,
\end{equation}
since $8.34/2=4.17$ and $.242/3<.081$.

The sign information now lets us estimate the average pair excess in
$x$.  Taking the marginal expectation in Equation~\eqref{loc:xpair}
and multiplying by $d$ gives the upper bound
\[
 -\mu\mu_j-2(\mu\eta_j-a\eta\mu_j)-2a\eta\eta_j
       +2.55(\mu\nu_j+a^2\nu\mu_j).
\]
The terms $-2\mu\eta_j+C\mu\nu_j$ suggest keeping their
combination together.  Define
\[
 K_j=C\nu_j-2\eta_j,\qquad G=2a\eta+C a^2\nu.
\]
Using $2.55=C+.05$, an exact rearrangement of the bound gives
\begin{equation}\label{loc:finalincrement}
 (-\mu+G)\mu_j+(\mu+a\eta)K_j-Ca\eta\nu_j
       +.05(\mu\nu_j+a^2\nu\mu_j).
\end{equation}
The term $-Ca\eta\nu_j$ is nonpositive.  The remaining new
quantity is the average of $K_j$, which we estimate before returning
to the $x$ excess.

The standalone normalized $K$ increment is $G$, with
\begin{equation}\label{loc:G}
 0\le G\le[2(1.417)+2.5(1.235)]\mu^2
       =5.9215\mu^2<5.922\mu^2.
\end{equation}
Equation~\eqref{loc:kpair} bounds its normalized excess by
\[
 (-2\mu-22a\eta+55.5a^2\nu)\mu_j
 +(22\mu+7a\eta)\eta_j+26\mu\nu_j.
\]
After discarding $-22a\eta\le0$, the coefficient of $\mu_j$ divided
by $\mu$ is at most
$-2+55.5(1.235)(.05)=1.427125<1.43$.
The coefficient of $\eta_j$ is positive, allowing use of
\eqref{loc:etaimproved}.

The average of a partial sum equals the weighted average of its
increments.  More precisely, in the regular case, for
$w_j=(d-1-j)/d$,
\[
 \operatorname{av}K_j
 =\operatorname{av}\{w_j d(K_{j+1}-K_j)\}.
\]
In the Poisson case Fubini gives the same identity with weight $1-s$
and derivative in $s$.  Direct finite sums, or the corresponding
integrals, give
\begin{equation}\label{loc:weightedaverages}
 \operatorname{av}w=\frac\alpha2,\quad
 \operatorname{av}(ws)\le\frac\alpha6,\quad
 \operatorname{av}(ws^2)\le\frac\alpha{12}.
\end{equation}
The latter two discrete values are
$(d-2)(d-1)/(6d^2)$ and
$(d-2)(d-1)^2/(12d^3)$, respectively.
Since $s^3\le s^2$, applying
\eqref{loc:partialmu}, \eqref{loc:partialnu},
\eqref{loc:etaimproved}, and~\eqref{loc:G} gives
\begin{align}
 \frac{\operatorname{av}K_j}{\alpha\mu^2}
 &\le5.922/2+1.43(1.14)/6\nonumber\\
 &\quad+(22+7(1.417)(.05))(.05)(4.17+.081)/12\nonumber\\
 &\quad+26(.05)(1.235/6+3.705/12)\nonumber\\
 &=\frac{20640565669}{4800000000}<4.31.
 \label{loc:Kaverage}
\end{align}
No nonnegativity of $K_j$ has been assumed.

We can now average Equation~\eqref{loc:finalincrement}.
The third term is nonpositive.  Use the lower bound on $\mu_j$ for
$-\mu\mu_j$ and the upper bound for $G\mu_j$, where $G\ge0$.
For the $K_j$ term first use~\eqref{loc:Kaverage}, since
$\mu+a\eta>0$, and then bound its positive coefficient.
After division by $\alpha\mu^2>0$, the average of
\eqref{loc:finalincrement} is at most
\begin{align*}
 &-(1/2-.05/3)+5.922(.05)(.57)
       +(1+1.417(.05))(.05)(4.31)\\
 &\hspace{2em}+.05(.05)(1.235/2+3.705/3+1.235(.57))
       =-\frac{2321911}{30000000}<0.
\end{align*}
This contradicts the zero actual excess average and proves the
proposition.  The omitted case $\mu=0$ is exactly the uniform posterior
experiment.
\end{proof}
